\chapter{Introducción}
\label{ch:ch1_intro}

% [BORRADOR v1] Material de partida: docs/plan.md, docs/arquitectura.md, docs/via_asic.md.

\section{Motivación: una clave vale lo que vale su entropía}
\label{sec:motivation}

Toda conexión segura que abre un dispositivo embebido, toda imagen de
firmware que autentica y toda credencial que almacena descansan sobre una
clave, y toda clave descansa sobre una fuente de aleatoriedad. Si esa
fuente es más débil de lo supuesto, ninguna sofisticación criptográfica
posterior recupera lo que se perdió al principio: una clave de 128 bits
extraída de un generador con 40 bits efectivos de entropía es una clave
de 40 bits. No es una preocupación teórica. La inyección de frecuencia en
la línea de alimentación de un microcontrolador seguro comercial redujo
su espacio de claves de $2^{64}$ a unas
3300~\cite{markettos_2009_frequency}; un generador certificado de tarjeta
inteligente resultó estar sesgado pese a superar las baterías estadísticas
estándar~\cite{hurleysmith_2018_certifiably}; y el generador hardware de
uno de los microcontroladores IoT más extendidos, el ESP32, devuelve en
silencio valores pseudoaleatorios cuando su radio está apagada sin dejar
de superar todos los tests de caja negra~\cite{blancoromero_2026_entropy},
algo que este trabajo mide directamente con la prueba de reinicios de
SP~800-90B (Sección~\ref{sec:esp32_90b}).

El hilo común es que la aleatoriedad se ha juzgado mirando la salida. Las
baterías estadísticas como NIST SP~800-22 comprueban si un flujo de bits
\emph{parece} uniforme e independiente; no pueden distinguir una fuente de
ruido física de un generador determinista cuya semilla es una constante, y
el NIST lo dijo formalmente en 2022 al anunciar la revisión de esa
publicación, ``rechazando su uso para evaluar generadores criptográficos
de números aleatorios''~\cite{nist_2022_decision}. Las normas actuales
--- AIS~20/31 en Europa~\cite{peter_2024_ais31} y SP~800-90B en
Estados Unidos~\cite{turan_2018_sp80090b} --- piden algo distinto: un
modelo estocástico de la fuente de ruido, una estimación de su
min-entropía derivada de ese modelo y tests de salud que vigilen la propia
fuente en lugar de su salida. Es un problema de microelectrónica antes que
de criptografía, y es el problema que aborda este TFM.

\section{Del estudio de caja negra del ESP32 a un motor de caja blanca}
\label{sec:background}

Este trabajo continúa un proyecto anterior del autor que caracterizó como
caja negra el generador hardware de números aleatorios del ESP32: se
ejercitó el mismo periférico en cuatro condiciones de funcionamiento, su
salida se pasó por una implementación autovalidada de los quince tests de
SP~800-22 y se propuso un esquema de recolección de claves que combinaba
varias fuentes internas del chip. Su hallazgo central fue metodológico:
los tests se superaron en condiciones en las que la entropía no podía ser
la que los tests daban a entender, que es exactamente el punto ciego
descrito antes. Dos de sus cifras también tuvieron que revisarse a
posteriori --- se informó de una correlación con la temperatura a partir
de un sensor cuya resolución era un solo escalón en todo el rango
observado, y se infirió una cifra de estabilidad donde podría haberse
medido --- y ambas lecciones se convierten en reglas en la
Sección~\ref{sec:conventions}.

El paso natural siguiente era construir el generador en lugar de medir el
de otro: un motor cuya fuente de entropía esté diseñada, cuyo jitter se
mida dentro del dispositivo, cuya min-entropía se deduzca de un modelo y
cuyas claves produzcan y usen primitivas verificadas frente a las normas.
El ESP32 cambia de papel: de objeto de estudio pasa a juez. Tiene sus
propios aceleradores AES y SHA y una biblioteca criptográfica madura, de
modo que puede recalcular todo lo que produce el motor con una
implementación que no comparte nada con él, a través de un bus real.

El motor está diseñado para silicio. Los elementos seguros --- pequeños
circuitos integrados que generan, almacenan y usan claves en nombre de un
microcontrolador anfitrión a través de I2C --- son una categoría de
producto madura; la \emph{transparencia} de su entropía no lo es. De tres
componentes de uso extendido, uno se vende con una validación NIST de su
fuente de entropía y dos no mencionan SP~800-90 en ningún lugar de su
documentación~\cite{microchip_atecc608b_ds}. La FPGA de este trabajo es,
por tanto, el vehículo de verificación, no el objetivo: permite ejecutar
millones de vectores de prueba contra el diseño en minutos, algo que
ningún simulador a nivel de transistor puede hacer, antes de que la misma
descripción se sintetice sobre celdas estándar.

\section{Objetivos}
\label{sec:objectives}

\begin{enumerate}
\item[O1] Diseñar a nivel de transistor una fuente de entropía de
oscilador en anillo, con habilitación y muestreador, y simularla con
ruido, declarando la resolución numérica del banco antes de leer de él
ninguna cifra.
\item[O2] Extraer el jitter de las formas de onda simuladas con un
estimador que separe las componentes de cuantificación, térmica y de
parpadeo, validado previamente sobre datos sintéticos de jitter conocido.
\item[O3] Convertir el jitter en una cota de min-entropía mediante un
modelo estocástico contrastado con un cálculo exacto, y deducir el
divisor del muestreador que cumple con margen el objetivo de AIS~20/31.
\item[O4] Escribir el motor digital --- tests de salud, acondicionamiento
basado en AES y CTR\_DRBG, AES-CMAC y una interfaz de registros I2C ---
como RTL independiente de la tecnología, verificado frente a vectores
oficiales.
\item[O5] Verificar funcionalmente el motor en una FPGA con el ESP32 como
oráculo independiente: pruebas de respuesta conocida, interoperabilidad
cruzada y reto-respuesta.
\item[O6] Realizar una campaña de entropía sobre silicio real (los anillos
de la FPGA) con los estimadores de SP~800-90B y comparar con el generador
del ESP32 bajo la misma metodología.
\item[O7] Sintetizar el motor sobre una biblioteca de celdas estándar
abierta e informar de su área en puertas equivalentes, con los factores
de crecimiento hasta el área de núcleo.
\item[O8] Comparar el diseño con el generador del ESP32, con elementos
seguros comerciales y con circuitos publicados, con incertidumbres
declaradas.
\end{enumerate}

El resultado mínimo defendible es O1--O3 junto con O5: una fuente de
entropía diseñada desde el transistor, con su jitter extraído y convertido
en una cota, y un motor digital verificado frente a una implementación
independiente. Un resultado negativo bien medido cierra el trabajo igual
que uno positivo: si el exponente de acumulación en hardware es dos en
lugar de uno, domina el ruido de parpadeo y el divisor no puede crecer
indefinidamente, y eso es un hallazgo.

\section{Alcance y exclusiones}
\label{sec:scope}

Cuatro cosas quedan deliberadamente fuera del alcance, cada una por una
razón explícita. \emph{Fabricación}: una lanzadera académica es asequible
y el motor cabe en el bloque más pequeño que se ofrece, pero el plazo
desde el tape-out hasta tener muestras encapsuladas y caracterizadas
excede lo que cabe en un TFM; el coste y el camino se documentan, no se
ejecutan. \emph{Criptografía asimétrica}: un multiplicador de curva
elíptica es un proyecto en sí mismo y consumiría todo el presupuesto de
área; el motor proporciona claves simétricas y autenticación.
\emph{Resistencia a canales laterales}: el modelo de amenaza es un
adversario lógico y remoto; las fugas de temporización se evitan por
construcción, pero no se intenta ningún enmascaramiento ni evaluación
basada en trazas, y se señala el propio laboratorio del instituto para
ese tipo de trabajo como la continuación natural. \emph{Certificación
formal}: los criterios de AIS~20/31 y SP~800-90B y los vectores oficiales
se usan en todo el trabajo, y no se reivindica ninguna validación.

\section{Metodología y convenciones de reproducibilidad}
\label{sec:conventions}

Cuatro convenciones se aplican a toda cifra de este documento. Primera,
cada número lleva la etiqueta \meas{}, \simu{} o \est{}, y va acompañado
de su incertidumbre o de la fuente de la que procede. Segunda,
\textbf{la resolución del instrumento se declara antes de interpretar
cualquier tendencia}, sea el instrumento un sensor, un estimador o un
simulador --- el suelo de ruido numérico del banco SPICE y la
cuantificación de un sensor de temperatura se tratan igual. Tercera, las
magnitudes se miden en lugar de inferirse siempre que exista una medida
directa. Cuarta, la novedad se expresa como ``no encontrado en la
literatura consultada'', nunca como ``no existe''; una afirmación de este
tipo en un borrador temprano de este trabajo tuvo que corregirse al
encontrarse un artículo de 2024 que medía exactamente la magnitud en
cuestión, y la corrección se deja visible en las notas de trabajo.

Todo se publica: el VHDL y sus bancos de pruebas, los modelos y el
análisis en Python, el firmware, el flujo de síntesis, los vectores de
prueba oficiales utilizados y los datos brutos de cada campaña, con las
órdenes que reproducen cada tabla.

\section{Estructura del documento}
\label{sec:structure}

El Capítulo~\ref{ch:ch2_soa} revisa los generadores basados en
osciladores en anillo, sus modelos estocásticos, las normas y el panorama
comercial. El Capítulo~\ref{ch:ch3_entropy} construye la cadena desde el
transistor hasta la min-entropía y valida cada eslabón. El
Capítulo~\ref{ch:ch4_engine} describe el motor digital y el
Capítulo~\ref{ch:ch5_verification}, su verificación. El
Capítulo~\ref{ch:ch6_asic} sintetiza el motor sobre celdas estándar y
expone el camino y el coste de un circuito integrado. El
Capítulo~\ref{ch:ch7_results} presenta la campaña en hardware y las
comparaciones, y el Capítulo~\ref{ch:ch8_concl} concluye.
